\vspace{10pt}\noindent\textbf{Thực nghiệm: Xếp chồng các lớp Hồi quy (Stacked Bidirectional RNN)} \hfill \\

Để tăng cường khả năng học các biểu diễn phức tạp và trừu tượng (Hierarchical Features), một kỹ thuật phổ biến trong Deep Learning là xếp chồng (Stack) nhiều lớp ẩn lên nhau. Ở thực nghiệm này, chúng ta sẽ xây dựng kiến trúc \textbf{Stacked Bidirectional RNN}, gồm 2 lớp SimpleRNN hai chiều nối tiếp nhau, đồng thời sử dụng lớp nhúng từ (Embedding Layer) tự học từ đầu thay vì mượn trọng số GloVe như bài trước.

\vspace{10pt}\noindent\textbf{1. Tiền xử lý Dữ liệu} \hfill \\

Tương tự như thực nghiệm trước, bộ dữ liệu IMDB gồm 50.000 mẫu hoàn toàn cân bằng (25.000 Tích cực, 25.000 Tiêu cực). Chúng ta áp dụng kỹ thuật loại bỏ Tạp âm miền (như các từ \textit{movie}, \textit{film}) và chuẩn hóa chiều dài câu ở mức 500 từ.

\begin{lstlisting}[style=pythoncode, caption={Thiết lập cấu hình và Load dữ liệu}]
import os
import pandas as pd
import numpy as np
import tensorflow as tf
from tensorflow.keras import layers, models
from tensorflow.keras.preprocessing.text import Tokenizer
from tensorflow.keras.preprocessing.sequence import pad_sequences
from sklearn.model_selection import train_test_split
from sklearn.metrics import classification_report, confusion_matrix

VOCAB_SIZE  = 20000
MAX_LEN     = 500
EMBED_DIM   = 128
EPOCHS      = 10
BATCH_SIZE  = 128
\end{lstlisting}

\begin{figure}[H]
    \centering
    \begin{minipage}{0.48\textwidth}
        \includegraphics[width=\linewidth]{Chương_3/results/imdb_stacked_rnn/01_eda_overview.png}
        \caption{Phân phối nhãn và độ dài văn bản}
    \end{minipage}
    \hfill
    \begin{minipage}{0.48\textwidth}
        
\includegraphics[width=\linewidth]{Chương_3/results/imdb_stacked_rnn/02_wordclouds.png}
        \caption{WordCloud sau khi loại bỏ Tạp âm miền}
    \end{minipage}
\end{figure}

\begin{lstlisting}[style=pythoncode, caption={Vector hóa văn bản với Tokenizer tự học}]
X_train, X_test, y_train, y_test = train_test_split(df['clean'].values, df['label'].values, test_size=0.2, random_state=42, stratify=df['label'].values)

tokenizer = Tokenizer(num_words=VOCAB_SIZE, oov_token='<OOV>')
tokenizer.fit_on_texts(X_train)

X_train_pad = pad_sequences(tokenizer.texts_to_sequences(X_train), maxlen=MAX_LEN, padding='post', truncating='post')
X_test_pad  = pad_sequences(tokenizer.texts_to_sequences(X_test), maxlen=MAX_LEN, padding='post', truncating='post')
\end{lstlisting}

\vspace{10pt}\noindent\textbf{2. Xây dựng Kiến trúc Stacked BiRNN} \hfill \\

Trong Keras, để xếp chồng các lớp RNN, lớp RNN đứng trước bắt buộc phải trả về toàn bộ chuỗi đầu ra ở mỗi bước thời gian để làm đầu vào cho lớp RNN tiếp theo. Điều này được thực hiện bằng cách đặt \texttt{return\_sequences=True}.

\begin{lstlisting}[style=pythoncode, caption={Khởi tạo Stacked Bidirectional SimpleRNN}]
model = models.Sequential([
    layers.Input(shape=(MAX_LEN,)),
    layers.Embedding(VOCAB_SIZE, EMBED_DIM, input_length=MAX_LEN, name='embedding'),
    layers.SpatialDropout1D(0.2),
    
    # Layer 1: Return sequences for the next RNN layer
    layers.Bidirectional(layers.SimpleRNN(128, return_sequences=True), name='bi_rnn_1'),
    layers.BatchNormalization(),
    layers.Dropout(0.3),
    
    # Layer 2: Final RNN layer
    layers.Bidirectional(layers.SimpleRNN(64), name='bi_rnn_2'),
    layers.Dropout(0.3),
    
    layers.Dense(64, activation='relu'),
    layers.Dense(1, activation='sigmoid', name='output')
], name='BiRNN_IMDB')

model.compile(optimizer=tf.keras.optimizers.Adam(1e-3), loss='binary_crossentropy', metrics=['accuracy'])
\end{lstlisting}

\begin{lstlisting}[style=consoleoutput, caption={Output: Số lượng tham số khổng lồ}]
Total params: 2,676,225
Trainable params: 2,675,713
Non-trainable params: 512
\end{lstlisting}
Với việc tự học toàn bộ lớp Embedding kích thước 128 chiều và 2 tầng RNN, mô hình sở hữu đến hơn 2.6 triệu tham số huấn luyện (lớn hơn nhiều so với 99 ngàn tham số thực học của BiGRU dùng GloVe).

\vspace{10pt}\noindent\textbf{3. Quá trình Huấn luyện và Sự Gục ngã (Vanishing Gradient)} \hfill \\

\begin{lstlisting}[style=pythoncode, caption={Huấn luyện mô hình với EarlyStopping}]
history = model.fit(
    X_train_pad, y_train, epochs=EPOCHS, batch_size=BATCH_SIZE, validation_split=0.15,
    callbacks=[
        tf.keras.callbacks.EarlyStopping(monitor='val_loss', patience=2, restore_best_weights=True),
        tf.keras.callbacks.ReduceLROnPlateau(factor=0.5, patience=2, verbose=1)
    ], verbose=1
)
\end{lstlisting}

\begin{figure}[H]
    \centering
    \includegraphics[width=0.8\linewidth]{Chương_3/results/imdb_stacked_rnn/03_training_history.png}
    \caption{Lịch sử Huấn luyện: Accuracy thấp và Validation Loss trồi sụt bất thường}
\end{figure}

\textbf{Đánh giá lịch sử huấn luyện:} 
Dù mô hình sâu hơn và có lượng tham số khổng lồ, hiệu suất lại \textbf{kém một cách thảm hại}. Độ chính xác trên tập Validation chỉ lẹt đẹt ở mức dưới 70\%. 
Nguyên nhân gốc rễ chính là hiện tượng Triệt tiêu Đạo hàm (Vanishing Gradient). Khi văn bản quá dài (500 từ), cộng thêm việc đạo hàm phải truyền ngược qua \textbf{nhiều tầng RNN xếp chồng} (chứ không phải 1 tầng), giá trị Gradient gần như bị nhân với 0 và biến mất hoàn toàn ở các trọng số đầu. Mô hình không thể học được sự phụ thuộc xa trong ngữ cảnh điện ảnh.

\vspace{10pt}\noindent\textbf{4. Đánh giá Confusion Matrix} \hfill \\

\begin{lstlisting}[style=pythoncode, caption={Đánh giá hiệu năng bằng Confusion Matrix}]
y_pred_prob = model.predict(X_test_pad, verbose=0).flatten()
y_pred      = (y_pred_prob >= 0.5).astype(int)

cm = confusion_matrix(y_test, y_pred)
cm_norm = cm.astype(float) / cm.sum(axis=1)[:, np.newaxis]

fig, axes = plt.subplots(1, 2, figsize=(13, 5))
sns.heatmap(cm, annot=True, fmt='d', cmap='Blues', xticklabels=['Negative', 'Positive'], yticklabels=['Negative', 'Positive'], ax=axes[0])
axes[0].set_title('Confusion Matrix (Raw Count)')

sns.heatmap(cm_norm, annot=True, fmt='.2f', cmap='YlOrRd', xticklabels=['Negative', 'Positive'], yticklabels=['Negative', 'Positive'], ax=axes[1])
axes[1].set_title('Confusion Matrix (Normalized)')

plt.savefig(f'{SAVE_DIR}/04_confusion_matrix.png'); plt.show()
\end{lstlisting}

\begin{figure}[H]
    \centering
    \includegraphics[width=0.9\linewidth]{Chương_3/results/imdb_stacked_rnn/04_confusion_matrix.png}
    \caption{Ma trận nhầm lẫn: Tỷ lệ đoán đúng chỉ tương đương với đoán mò có tính toán}
\end{figure}

\begin{lstlisting}[style=consoleoutput, caption={Output: Báo cáo hiệu năng Precision và Recall}]
              precision    recall  f1-score   support

    Negative       0.68      0.67      0.68      5000
    Positive       0.68      0.68      0.68      5000

    accuracy                           0.68     10000
\end{lstlisting}

\textbf{Kết luận:}
Dù có Precision và Recall cân bằng (68\%), nhưng mức Accuracy tổng thể \textbf{68\%} là một bước lùi quá lớn so với con số 89\% của kiến trúc GRU kết hợp GloVe. 
Thực nghiệm này mang lại một bài học đắt giá trong Deep Learning: \textbf{Mạng nơ-ron không phải cứ càng Sâu (Deep/Stacked) và càng Nhiều tham số thì càng mạnh}. Nếu không giải quyết được căn bệnh Vanishing Gradient bằng các cấu trúc Cổng (Gates như LSTM, GRU), việc xếp chồng các lớp SimpleRNN chỉ khiến mô hình trở nên "mất trí nhớ" và hỗn loạn hơn khi xử lý chuỗi thời gian dài.
